Interfaces gráficas mediam a maioria das nossas interações com computadores,
seja em celulares e \emph{notebooks}, seja no navegador web.
\todo{Precisa de fonte para \enquote{a maioria}.}
Para \textcite[p. 73]{myers1994}, projetá-las e implementá-las é difícil por
natureza e continuaria sendo.
Cada clique, tecla ou toque do usuário gera um evento.
Programar interfaces exige coordenar esses eventos com o estado do programa e
com o que a tela mostra.
Essa coordenação é o tema deste trabalho.
Uma das diferenças da programação de interfaces para a dos outros programas
é o código escrito do avesso (\emph{inside-out}), dividido em sub-rotinas que o
\emph{toolkit} chama quando ocorrem ações do usuário \cite[p. 79]{myers1994}.
Outra é o estado da aplicação, que determina a aparência e o comportamento da
tela \cite[p. 230]{oney2012}.
Além disso, ao escolher uma tecnologia, o programador escolhe
também a forma de escrever a interface, e essa forma afeta a compreensão e a manutenção
do programa \cite[p. 134]{fischer2007}.

O recurso predominante para propagar eventos é o \emph{callback}, também chamado
de \emph{listener}: um bloco de código que roda sempre que um dado evento ocorre,
como o clique em um botão \cite[p. 7]{blackheath2016}.
Registrá-lo num produtor, que o chama de volta a cada ocorrência, é o
\emph{Observer Pattern} \cite[p. 7]{blackheath2016}.
O \emph{callback} é síncrono quando a função que o recebe o chama antes de
retornar, e assíncrono quando sua execução fica para algum tempo depois de a
função retornar \cite[p. 2]{gallaba2015}.
O navegador executa a aplicação numa \emph{thread} só, e, para que ela continue
respondendo ao usuário, os programadores recorrem às operações assíncronas
\cite[p. 368-369]{alimadadi2014}.
O \emph{callback} de um clique é desse tipo: registrado quando a tela é montada,
só roda quando o usuário clica, e os eventos do \gls{dom} estão entre as
operações assíncronas do navegador \cite[p. 5]{gallaba2015}.

\emph{Threads} exprimem sequência, e eventos exprimem dependência, e muitos
problemas vêm de escrever uma com o recurso da outra
\cite[p. 16]{blackheath2016}.
O \emph{callback} fragmenta o fluxo de controle: a sequência de passos que o
programador tem em mente para uma operação fica dividida entre vários
\emph{callbacks}, o que dificulta a manutenção e a compreensão do programa
\cite[p. 134-135]{fischer2007}.
Como não se sabe em que ordem os \emph{callbacks} serão chamados, coordenar as
alterações ao estado compartilhado pode ser desconcertante e está longe de ser
modular \cite[p. 926]{edwards2009}.
Essas alterações são efeitos colaterais: as modificações que um comando, como a
atribuição, faz no estado do programa \cite[p. 361]{hudak1989}.
Os \emph{callbacks} que alteram o \gls{dom} por efeitos colaterais muitas vezes
levam a um código interdependente, opaco e propenso a erros, o chamado código
espaguete \cite[p. 229]{oney2012}.
Em JavaScript, os \emph{callbacks} são frequentes: num estudo de 138 programas, em
média uma em cada dez definições de função recebe um \emph{callback}, e a maioria dos
\emph{callbacks} é aninhada \cite[p. 1]{gallaba2015}.

O JavaScript ganhou alternativas ao \emph{callback}: as \emph{promises}, que encadeiam
computações assíncronas \cite[p. 86:1]{madsen2017}, e as palavras-chave
\emph{async} e \emph{await}, que permitem escrever operações assíncronas em estilo
sequencial \cite[p. 9]{grolaux2026}.
As \emph{promises} achatam o aninhamento dos \emph{callbacks}, mas a lógica complexa
continua presa em cadeias de \emph{callbacks} \cite[p. 9]{grolaux2026}.
Numa implementação em JavaScript que guarda o estado da aplicação em variáveis
globais, o \emph{async} e o \emph{await} não ajudam a coordenar esse estado
\cite[p. 535 e 538]{berry2020}.

Outra alternativa ao \emph{callback} é a \gls{pr}.
Segundo \textcite[p. 1125-1126]{salvaneschi2017}, nela o programador
declara as relações entre os valores, e o ambiente de execução faz as
atualizações que elas exigem.
Assim, o programador não precisa cuidar da ordem dos eventos nem das
dependências entre as computações \cite[p. 2]{bainomugisha2013}.
A \gls{pr} serve às aplicações orientadas a eventos e interativas e leva às
linguagens de programação o modelo das planilhas eletrônicas, em que mudar
o valor de uma célula recalcula as fórmulas que dependem dela
\cite[p. 1-2]{bainomugisha2013}.

\textcite[p. 1125-1134]{salvaneschi2017} testaram uma afirmação repetida,
ainda sem evidência empírica suficiente, de que a \gls{pr} torna os programas
mais fáceis de compreender que o \emph{Observer Pattern}, o padrão tradicional da
\gls{poo} para aplicações interativas.
Num experimento controlado com 127 estudantes, quem leu os programas escritos
com \gls{pr} acertou mais perguntas sobre eles do que quem os leu com o
\emph{Observer Pattern} (médias de 8,13 e 7,05 acertos em dez, com diferença
significativa).
Também respondeu mais rápido em seis das dez tarefas, sem diferença
significativa nas outras quatro.

O código sequencial tem o problema oposto ao do \emph{callback}: a linguagem
obriga a fixar a ordem de execução mesmo quando ela não importa, e quem lê
precisa supor que importa \cite[p. 8-9]{moseley2006}.
Para \textcite[p. 1]{moseley2006}, essa preocupação com o fluxo de controle
está, ao lado do volume de código e do tratamento do estado, entre as
principais causas da complexidade dos sistemas de grande porte.
\textcite[p. 929]{edwards2009} sugere que boa parte do sequenciamento manual
das ações que a programação imperativa exige pode ser feita automaticamente,
ao menos nas aplicações interativas.
A programação declarativa expressa o que se computa, e não como, embora isso
seja questão de grau, e suas linguagens não têm estado implícito
\cite[p. 361]{hudak1989}.
Nas interfaces, o componente do React é uma função que especifica a tela de
forma declarativa e que, chamada com os mesmos argumentos, deve exibir a
mesma interface \cite[289:1]{lee2025}.

Este trabalho compara \emph{modelos de programação}: as técnicas e os
princípios de projeto que um modelo de computação torna possíveis
\cite[p. xiii e 29]{roy2004}.
\todo{O que se compara passa a se chamar modelo de programação, e \enquote{notação} fica no sentido das DCs, a forma escrita de cada tecnologia. Solid e Angular com signals são um modelo de programação com duas notações, que variam na montagem da tela, no modelo de componente e na API.}
Na programação de interfaces gráficas, o modelo de programação é o modo
de coordenar evento, estado e tela que uma tecnologia impõe.
Ele não é o paradigma, noção imprecisa que o modelo de computação torna
precisa \cite[p. xiii]{roy2004}.
Também não é a biblioteca nem o \emph{framework}, que nomeiam o produto: o
jQuery se diz biblioteca \cite{openjs2026}, o Web Component é recurso do
próprio navegador \cite{mozilla2026}, e os dois têm o mesmo modelo de
programação.
A \emph{notação} é a forma escrita de um modelo de programação numa
tecnologia: as palavras e os sinais com que o programador escreve, no
código-fonte da aplicação, o estado, os valores derivados dele e a tela
que os mostra.
Nas tecnologias que este trabalho compara, há três modelos de
programação, dois deles declarativos.
No \emph{imperativo com callbacks}, o programador escreve \emph{callbacks} que
mudam o estado e a tela por efeitos colaterais
\cites[][ p. 93]{disch2025}[][ p. 229]{oney2012}, como no jQuery
\cite{openjs2026} e nos Web Components sem biblioteca \cite{mozilla2026}.

Nos declarativos, o programador descreve a tela em função do estado, e o
ambiente de execução a mantém atualizada
\cites{solidjs2026}[][ p. 32]{sperber2025}.
No \emph{declarativo por re-renderização}, cada mudança de estado reexecuta o
componente, como no React \cite[p. 289:2 e 289:6]{lee2025}.
No \emph{declarativo por atualização granular}, cada \emph{signal} avisa os valores
e trechos da tela que dependem dele, como no Solid e no Angular com
\emph{signals} \cite{solidjs2026,google2026}.
No Solid, só esses são recalculados \cite{solidjs2026}; no Angular, quando o
\emph{signal} muda, o \emph{template} do componente que o lê volta a ser renderizado
em algum momento \cite{google2026}.
O modelo de programação por atualização granular usa as abstrações da \gls{pr}, os
valores que variam no tempo, chamados \emph{signals}, cujas alterações se propagam
automaticamente \cite[p. 953]{salvaneschi2015}.
O Angular acrescentou essa abstração à sua \gls{api} \cite[p. 18]{lima2024},
e o Solid e o Angular com \emph{signals} aparecem entre as implementações da
\gls{pr} \cite[p. 57]{tsukanova2026}.
Cada modelo de programação aplica conceitos de um ou mais paradigmas.
Em quase todo problema, cada parte pede o seu conjunto de conceitos
\cite[p. 10]{roy2009}, e as implementações deste trabalho também os
misturam.

As pesquisas de uso medem tecnologias, não modelos de programação, e por
isso o uso de cada modelo de programação se afere pelas tecnologias que o
adotam.
Na questão sobre tecnologias web do Stack Overflow Developer Survey 2025,
React, jQuery e Angular estão entre as sete mais usadas, com 44,7\%, 23,4\% e
18,2\% dos que a responderam \cite{stackoverflow2025}.
Assim, o uso do declarativo por re-renderização se afere pelo React, e o
do imperativo, pelo jQuery, que também está em 65,5\% dos
\emph{sites} que o W3Techs monitora \cite{w3techs2026}.
O do declarativo por atualização granular se afere pelo Angular, embora o
Stack Overflow não separe quem o usa com \emph{signals}, e pelo Solid, usado por 10\%
dos respondentes do State of JS 2025 \cite{devographics2026}.
Nenhuma dessas pesquisas mede o uso de Web Components sem biblioteca.
Eles entram por ser a forma do próprio navegador de criar elementos
personalizados \cite{mozilla2026}.

Para discutir o quanto uma notação facilita a compreensão e a
manutenção, há as \glspl{dc}, propostas por \textcite{green1989}.
Elas são um vocabulário para descrever a usabilidade de notações e de outros
artefatos de informação \cite[p. 327]{blackwell2001}.
Trabalhos anteriores compararam por elas paradigmas e linguagens de
programação de interfaces e bibliotecas de \gls{pr}.
\textcite{mernik2009} compararam, com 36 estudantes de graduação, uma linguagem de
domínio específico para descrever interfaces gráficas, o \gls{xaml}, e uma biblioteca de
aplicação, o Windows Forms em C\#.
\textcite[p. v, 11 e 104]{kiss2014} comparou a \gls{poo} e a \gls{pf}
em sete tarefas de interface, com JavaFX e ScalaFX, e em parte delas com
Scala.Rx, ReactFX e Elm, e deixou como trabalho futuro comparar os
\emph{frameworks} de interface do navegador.
\textcite[p. 1510 e 1513]{zimmerle2025} avaliaram com usuários duas
bibliotecas de \gls{pr} para JavaScript, RxJS e Bacon.js, em tarefas de
requisições \gls{http}.
Fora das \glspl{dc}, \textcite[p. 15]{grolaux2026} compararam
conceitualmente, numa prova de conceito em JavaScript, o async/await com
\emph{listeners}, barramento de eventos, \emph{streams} funcionais e \gls{pr} no tratamento
dos eventos de interface.

Nenhum desses trabalhos compara, sobre as mesmas tarefas, os três modelos
de programação nas tecnologias web mais usadas.
Os trabalhos encontrados na revisão bibliográfica também não avaliam pelas
\glspl{dc} a notação do React nem a do Solid e do Angular com \emph{signals}.
A lacuna se apoia nessa revisão, feita por palavras-chave e por
\emph{snowballing} \cite[p. 3-4]{wohlin2014}, sobretudo para a frente, em duas
rodadas, em setembro e outubro de 2026, e descrita no Apêndice \ref{apend:2}.
A revisão não viu os trabalhos que citam \textcite{kiss2014}, que a base
usada não indexa, e a lacuna vale para o que ela encontrou.
Pergunta-se, então: pelas \glspl{dc}, como difere a usabilidade da
notação com que se programam interfaces gráficas em cada tecnologia entre
os três modelos de programação~--- o imperativo com \emph{callbacks}, o
declarativo por re-renderização e o declarativo por atualização granular?
Conhecer essas diferenças pode ajudar quem programa interfaces a escolher uma tecnologia
sabendo o que o modelo de programação dela facilita e o que dificulta em
cada \emph{problema de coordenação}.
Dentro de um mesmo modelo de programação, a notação também muda de uma
tecnologia para outra, como entre o Solid e o Angular com \emph{signals}, e a
avaliação registra quando essa diferença muda a conclusão.
Um problema de coordenação é uma exigência que a interface faz ao modelo de
programação.

O objetivo geral é comparar, segundo as \glspl{dc}, a usabilidade das
notações dos três modelos de programação em interfaces típicas da web,
em TypeScript, e em parte delas no Android, em Kotlin.
Especificamente:

\begin{itemize}
\item implementar cinco tarefas de interface (Contador, Formulário com validação,
Busca com sugestões, Lista filtrável e Carrinho) com Web Component,
jQuery, React, Solid e Angular com \emph{signals}, e o Contador, o Formulário e
a Lista no Android, com Views e Jetpack Compose;
\item avaliar as implementações por oito \glspl{dc};
\item sintetizar, por problema de coordenação, o que cada modelo de
programação facilita e o que dificulta, como no estado derivado e na
assincronia.
\end{itemize}

\todo[inline]{Saiu o objetivo de demonstrar a PF com processamento de listas, herdado de 2017: fundamentar conceitos cabe à seção 2, e não é passo para o objetivo geral. Os outros mudam com a pergunta, que deixou de comparar paradigmas e passou a comparar três modelos de programação pela usabilidade das notações, segundo as DCs.}

Quanto aos objetivos, a pesquisa é \emph{exploratória}: busca compreensão nova
e gera ideias e hipóteses para outras pesquisas
\cite[p. 135]{runeson2009}.
A pergunta de pesquisa, como a usabilidade difere entre os modelos de
programação, é do tipo que \textcite{easterbrook2008} chamam de
descritivo-comparativo, das perguntas sobre como uma coisa difere de outra,
e põem entre as perguntas exploratórias.
As hipóteses desta pesquisa servem a experimentos controlados como o de
\textcite{salvaneschi2017}.
\todo[inline]{Em relação ao projeto de 2017: a pergunta deixa de comparar paradigmas, PF e PR frente à POO com callbacks, e passa a comparar três modelos de programação pelas notações, com as DCs como critério explícito. A demonstração de PR e callbacks dá lugar a cinco tarefas de interface, implementadas em cinco tecnologias web, e três delas também no Android. O JavaScript passa a TypeScript, e entra o Kotlin. A classificação quanto aos meios também muda (nota adiante). A natureza aplicada sai, e a classificação exploratória passa a se apoiar em Runeson e Höst (2009) e em Easterbrook et al. (2008), no lugar de Leal (2011).}

Quanto aos meios, o trabalho é uma avaliação qualitativa, pelas \glspl{dc}, de
implementações de um mesmo conjunto de tarefas.
O conjunto de tarefas é um \emph{benchmark} de notação: testes que representam
as tarefas da prática e servem para comparar técnicas alternativas por
medidas de desempenho, que podem ser qualitativas e feitas por uma pessoa
\cite[p. 76]{sim2003}.
A medida, neste trabalho, é a avaliação pelas \glspl{dc}, como nos
\emph{benchmarks} de resultado qualitativo do projeto de linguagens
\cite[p. 35]{erdweg2015}, entre os quais está o 7GUIs
\cite[8:20]{lu2021}.
Aqui, as técnicas são os três modelos de programação, e as tarefas, cinco interfaces
típicas da web, que exercitam cinco problemas de coordenação:

\begin{description}
\item[{Contador}] o caminho mais simples, do evento ao estado e à tela;
\item[{Formulário com validação}] o estado derivado e os campos que dependem
uns dos outros;
\item[{Busca com sugestões}] a assincronia (espera, respostas fora de ordem e
cancelamento);
\item[{Lista filtrável}] a lista derivada por busca, filtro e ordenação, além
da montagem dos itens na tela;
\item[{Carrinho}] o estado compartilhado entre componentes.
\end{description}

O Contador, o Formulário e a Lista adaptam o \emph{Counter}, o \emph{Flight Booker} e
o \emph{CRUD}, três das sete tarefas do 7GUIs, de \textcite[p. 11, 17, 25 e 34]{kiss2014}.
A Busca e o Carrinho não têm equivalente no 7GUIs.
\todo[inline]{O projeto classificava o trabalho, quanto aos meios, como estudo de casos múltiplos (Leal 2011; Yin 2001). As definições reunidas por Runeson e Höst (2009, p. 134 e 139) pedem um fenômeno contemporâneo no seu contexto real e excluem os programas de exemplo. Por isso a classificação saiu, e as interfaces implementadas passaram a se chamar tarefas, como no 7GUIs (Kiss 2014) e em Sim et al. (2003). Registrar essa mudança no texto?}

Na web, o modelo de programação imperativo com \emph{callbacks} entra com o Web
Component e o jQuery, e o delineamento compara dois pares de tecnologias próximas.
React e Solid montam a tela em JSX, marcação parecida com \gls{html}
escrita dentro do JavaScript \cite{meta2026c}.
O que os separa, sobretudo, é o modo de coordenar: a cada mudança de
estado, o React reexecuta o componente inteiro, e o Solid atualiza só os
trechos da tela que dependem do \emph{signal} alterado
\cites[][ 289:38]{lee2025}{solidjs2026}.

Solid e Angular com \emph{signals}, ao contrário, coordenam do mesmo modo: os
dois guardam o estado em \emph{signals} e derivam valores de outros \emph{signals}
\cite{solidjs2026a,google2026}.
Por isso ficam no mesmo modelo de programação, e o par mostra como a
notação varia entre tecnologias em três pontos.
O Solid monta a tela em JSX numa função, e o Angular, num \emph{template}
baseado em \gls{html}, com construtos próprios \cite{google2026a}.
O Angular escreve o componente como classe, com decorador
\cite{google2026b}, e o componente participa da injeção de dependências
\cite{google2026c}.
E as \glspl{api} dos dois diferem no que vai além do \emph{signal} e do valor
derivado, como o cancelamento das requisições na Busca.

No Android, o Formulário e a Lista são implementados com Views e Compose,
e a comparação dos dois diz se as conclusões tiradas na web se repetem em outra plataforma e outra
linguagem.
As Views ficam no imperativo com \emph{callbacks}, e o Compose, no declarativo
por re-renderização, porque chama de novo as funções da tela quando o estado
muda \cite{android2026}.
O Contador também é implementado no Android, como introdução à sintaxe,
sem análise própria.

O delineamento mantém fixo o que não pertence ao modelo de programação nem à
notação.
A linguagem é uma só em cada plataforma: na web, o \emph{TypeScript}, que o
Angular exige \cite{google2026b} e as demais tecnologias também usam, e, no
Android, o Kotlin.
As implementações de uma tarefa seguem a mesma especificação, com os mesmos
textos, regras e dados, e cada uma usa só o que a própria tecnologia traz,
sem bibliotecas de formulário.
As regras que não dependem da interface, como o formato das datas, a
validação do e-mail e o catálogo de produtos, ficam num \emph{domínio} comum
às implementações de cada plataforma, em TypeScript na web e em Kotlin no
Android.

Uma \emph{rotina} de teste por tarefa executa as verificações da especificação
contra cada implementação; na web, ela roda no console do navegador.
\emph{Capturas} de tela dos estados escolhidos de cada tarefa, geradas
automaticamente, denunciam desvios de comportamento ou de aparência,
porque devem sair idênticas entre as tecnologias web.
As implementações, as rotinas e as capturas ficam num repositório público,
em \url{https://github.com/iquabius/unemat-tcc}.
Nele, as versões de linguagens, bibliotecas e ferramentas ficam
fixadas\footnote{No \texttt{package-lock.json} e no \texttt{libs.versions.toml}.
Na web, TypeScript 6.0.3, jQuery 4.0.0, React 19.3.0, Solid 1.9.15 e Angular
22.2.0, com Vite 8.3.1 para o \emph{build} (no Angular, o \texttt{@angular/build} 22.2.0)
e Playwright 1.63.0 para as capturas.
No Android, Kotlin 2.4.20, Jetpack Compose pela \gls{bom}
2026.09.00 e, para as capturas, Robolectric 4.17 com Roborazzi 1.75.0.}, para que as implementações possam ser reproduzidas.

A avaliação cabe ao autor, que examina o código das implementações dimensão
a dimensão, sem participantes.
É uma avaliação analítica, feita por especialista, sem a participação de
usuários \cite{blandford2008}, como a de \textcite[p. 8]{kiss2014} pelas
\glspl{dc}.
Ela serve de alternativa aos estudos com programadores, que comparam
linguagens a custo alto e com resultados muitas vezes inconclusivos
\cite[p. 9]{sadowski2011}, e as \glspl{dc} pedem menos tempo que eles
\cite{dagit2006}.

A avaliação usa oito das catorze \glspl{dc} da lista de
\textcite[p. 12]{kiss2014}.
Seis são as que ele usou \cite[p. 12-14]{kiss2014}: \emph{nível de
abstração}, \emph{proximidade de descrição}, \emph{dependências ocultas},
\emph{propensão a erros}, \emph{concisão} e \emph{viscosidade}.
As outras duas, \emph{expressividade} e \emph{operações mentais difíceis}, vêm das
cinco que \textcite[p. 15]{kiss2014} considerou e deixou de fora por terem
pesado pouco.
A seu ver, elas provavelmente teriam sido úteis se as linguagens e os
\emph{toolkits} comparados fossem muito mais diferentes~--- aqui, cada
tecnologia traz as próprias abstrações.
Não entram as outras três dessas cinco, pelas razões discutidas
na seção \ref{chap:cases}.
Também ficam de fora, como em \textcite[p. 15]{kiss2014}, as três que
tratam das ferramentas e do processo de escrita, a visibilidade, a
análise progressiva e a provisoriedade, porque o objeto da avaliação é o
código pronto.

A seção \ref{chap:results} compara os modelos de programação a partir das
conclusões da avaliação e abre com uma tabela-síntese, com uma linha por
\gls{dc} e uma coluna por modelo de programação.
Cada célula resume a conclusão, indica de que tarefa vem a evidência e
registra quando as duas tecnologias da coluna divergem, o Solid e o
Angular com \emph{signals} ou o Web Component e o jQuery.
No primeiro par, ela diz também se a diferença vem da montagem da tela,
do modelo de componente ou da \gls{api}.
Como cada tarefa exercita um problema de coordenação, a síntese por
problema reúne as células que vêm da tarefa dele.
Segue uma seção por \gls{dc}, com os trechos de uma tarefa escolhida,
antes da avaliação, pelo problema de coordenação que a dimensão mais toca.
Em cada seção, a análise aponta se a diferença está na coordenação, na
montagem da tela ou na ligação entre as duas.
Nas seções em que o Formulário ou a Lista mostram diferença entre Views e
Compose, a análise diz se a conclusão da web se repete no Android.
A concisão tem ainda o apoio de uma contagem das linhas de cada
implementação.
\todo[inline]{Termos: \enquote{comparar} fica para o objetivo, \enquote{avaliação} para a aplicação das DCs ao código e \enquote{análise} para a parte que apresenta os resultados. Considerar um termo só, \enquote{avaliação}, para a atividade toda. A favor: uma palavra por conceito, sem o leitor se perguntar se são três atividades. Contra: mudam mais frases, e o objetivo geral perde o verbo \enquote{comparar}, que diz o que se faz entre os modelos de programação.}

O delineamento tem limitações em três pontos: quem avalia, o instrumento e o
objeto.
Como a avaliação cabe ao autor, nenhum outro avaliador confere as conclusões
em cada dimensão.
Em métodos de inspeção, avaliadores diferentes acham problemas bem
diferentes na mesma interface \cite[p. 183]{hertzum2003}.
Além disso, quem avalia pelas \glspl{dc} o que construiu tem conflito de
interesse e tende a não ver os problemas que se acostumou a contornar
\cite{dagit2006}.
A familiaridade do autor com as tecnologias é desigual, maior com o React.
Na avaliação, que é subjetiva, esse desnível pode favorecer a notação do
React.

Por ser analítica, e não um experimento, a avaliação pelas \glspl{dc} aponta
diferenças que ainda precisam ser testadas com programadores.
O trabalho escolheu as oito \glspl{dc} antes da avaliação, e escolher um
subconjunto pela relevância para a atividade pode deixar de fora
propriedades importantes da notação \cite[p. 265]{britton2001}.
Para \textcite[p. 334]{green2006}, duas delas, a expressividade e as
operações mentais difíceis, sempre foram mal descritas.
Das três que não entram, a exclusão mais discutível é a do compromisso
prematuro, que também aparece na própria notação, quando um trecho escrito
numa construção pede outra e mudar custa caro \cite[p. 155 e 157]{green1996d}.

As tarefas são pequenas, e o que só aparece em programas de grande porte ou ao
longo da manutenção fica de fora.
O autor definiu o conjunto de tarefas para este trabalho.
Um conjunto criado por um só pesquisador e pouco usado tende a ter menos
impacto que um adotado por uma comunidade \cite[p. 76]{sim2003}.
O delineamento compara categorias, e não artefatos: cada modelo de programação
tem uma ou duas tecnologias na web, e a conclusão sobre ele generaliza a
partir delas.
As \glspl{dc} também dizem pouco das escolhas de alto nível, como o
modelo de computação, e a relevância delas começa quando se definem os
detalhes \cite[p. 139]{green1996d}.
Por isso, a avaliação compara os modelos de programação pelos detalhes que
a notação de cada tecnologia exige no código-fonte.
Na web, o imperativo com \emph{callbacks} difere dos declarativos em mais de um
ponto ao mesmo tempo, e nenhum par isola o modelo de programação dele.
No Android, Views e Compose são o par mais próximo, mas mudam também a
montagem da tela, e o Android confere só o imperativo com \emph{callbacks} e
o declarativo por re-renderização.

Além desta introdução, o texto tem quatro seções.
A seção \ref{chap:prog} reúne os conceitos de programação em que o
trabalho se apoia, da programação de interfaces gráficas à \gls{pf} e à \gls{pr}.
A seção \ref{chap:cases} apresenta as \glspl{dc} da avaliação e descreve as
implementações das cinco tarefas.
A seção \ref{chap:results} compara os três modelos de programação pelas
\glspl{dc}, e a
seção \ref{chap:conclusion} responde à pergunta de pesquisa.
\todo{Os capítulos ainda seguem o projeto de 2017 e mudam com a seção de trabalhos relacionados; este parágrafo se revê no fechamento.}
